\section{Discussion}

\subsection{Principal Findings}

This study formulates tissue-specific MHC-I presentation preference as a paired multi-task learning problem. Rather than asking whether a peptide can bind to or be presented by an MHC molecule in general, the benchmark asks which of two peptides with matched MHC and parent-protein evidence is preferentially associated with a particular tissue--MHC task. This formulation controls two major sources of trivial discrimination and focuses evaluation on relative tissue-conditioned evidence.

Seven findings summarize the study. The paired task is learnable in the 157-task human benchmark, where the TissuePMHC human model reaches mean task AUROC 0.8448. Global and HLA-specific branches contribute complementary rankings on the standard benchmark. A frozen Factorized MMoE reaches AUROC 0.8562 across 24 mouse tasks, replicating task learnability under the related construction. Matched peptide-disjoint OOF reduces AUROC to 0.7652 in human and 0.7529 in mouse, demonstrating an unseen-peptide generalization gap. Strict architecture controls retain a clear human auxiliary-supervision benefit but do not establish superiority of the full human model over the auxiliary dual branch or of Factorized MMoE over simpler mouse controls. Capacity-matched MHC-only controls remain lower under every protocol, including strict AUROCs of 0.6630 in human and 0.6412 in mouse. Finally, frozen general-pMHC predictors reach moderate AUROC, and adding their scores changes AUROC by at most 0.0041.

The standard benchmarks support comparisons within the original closed-task construction, whereas peptide-disjoint OOF tests held-out peptide entities within those tasks. These results are complementary; the full exclusions and deployment limits are consolidated in the Limitations section.

\subsection{Architecture Effects under Peptide Separation}

The standard-benchmark architecture advantage only partially persists after exact peptide identities are separated. In human, TissuePMHC remains above traditional peptide baselines, shared heads, and a plain dual branch on the strict folds. However, it is statistically comparable to the auxiliary dual-branch MLP. The auxiliary model itself exceeds both shared heads and the plain dual branch, making auxiliary supervision the most reproducible retained component benefit. These results do not support attributing the strict gain specifically to the multi-kernel encoder or to the dual/rank-fusion structure without auxiliary objectives.

The mouse interpretation is more conservative. Factorized MMoE, shared heads, and the BLOSUM62 random forest have comparable strict ensemble performance. Moreover, Factorized MMoE has the lowest mean single-seed AUROC of the three and the largest ensemble gain. The final mouse score therefore demonstrates learnability and a strong benefit from seed averaging, not a universal architecture advantage. More generally, architecture conclusions from the easier standard benchmark should not be assumed to survive entity separation without matched retraining on the same component folds.

\subsection{Relationship to General Peptide--MHC Prediction}

The TissuePMHC human model is not a replacement for a general peptide--MHC binding or presentation predictor. Conventional predictors estimate whether a peptide is compatible with an MHC molecule or likely to be presented under broadly pooled conditions. The present model conditions on a tissue--MHC task and ranks peptides within a matched evidence structure. Consistent with this relationship, MHCflurry presentation and NetMHCpan EL attain task-macro AUROCs of approximately 0.65--0.69: generic pMHC propensity clearly contributes to the benchmark, but remains below the tissue-conditioned models. The lower capacity-matched MHC-only results provide a more direct within-family control and show that removing tissue identity reduces performance. NetMHCpan EL is also stronger than binding-only BA, especially in human, indicating that the broader presentation signal is a more relevant external control than affinity alone.

This distinction also limits causal interpretation. Tissue-associated sequence patterns may reflect multiple processes, including source-protein expression, antigen processing, cell composition, experimental sampling, and residual study structure. Because the current model receives only peptide sequence and task identity, it cannot identify which upstream process produced an association. Cross-fitted stacking shows that frozen general-pMHC scores add at most 0.0041 AUROC to the current model, suggesting that most usable generic signal is already reflected in its ranking. This small increment should not be interpreted as a new model improvement or as evidence that the two score types are fully redundant.

\subsection{Potential Applications}

The benchmark may support several research applications. Tissue-conditioned scores could help prioritize candidate antigens whose presentation evidence is concentrated in a target tissue, rank potential normal-tissue off-targets during early therapeutic screening, and select peptides for follow-up immunopeptidomics experiments. The same framework could also contribute to candidate triage for T-cell receptor therapies, vaccines, and other immunotherapies when used alongside binding affinity, immunogenicity, expression, safety, and population-coverage evidence.

These are potential uses rather than validated clinical applications. The balanced evaluation tasks do not reproduce clinical prevalence, and retrospective database associations cannot establish safety or therapeutic benefit. Deployment would require prospective validation, calibrated decision thresholds, uncertainty estimates, and cohorts representative of the intended use. In particular, a low predicted score must not be treated as evidence that a peptide is absent from a normal tissue.

\subsection{Limitations}

The benchmark has several important limitations. The negative member of a pair is a pseudo-negative constructed from available evidence, not an experimentally confirmed non-presented peptide. Incomplete tissue sampling and limited detection depth can therefore convert unobserved positives into apparent negatives. This issue is especially relevant for peptides reported across multiple tissues and for tissues with heterogeneous cellular composition.

The source data are observational and inherit biases from IEDB studies. Tissue coverage, HLA frequency, sample preparation, mass-spectrometry platform, reporting practice, and detection depth are uneven. Matching on MHC and parent protein reduces specific confounders but does not remove study, laboratory, donor, or batch structure. The deliberately balanced 1:1 test tasks are useful for controlled ranking comparisons but do not represent real-world prevalence; consequently, benchmark AUPRC and score thresholds should not be transferred directly to deployment settings.

Source-tissue strings are retained as exact nonempty IEDB labels. Synonyms, anatomical subregions, and differently formatted labels are not ontology-normalized, so semantically related tissues may be split into separate tasks. Resolving those labels would change task definitions and require rebuilding the benchmark and retraining all models; the present revision therefore documents the issue rather than mixing incompatible results.

The present work covers unmodified canonical 9-mer MHC-I peptides. It does not address variable-length MHC-I ligands, post-translationally modified peptides, spliced peptides, or MHC-II presentation. It also excludes explicit tissue transcriptomics, proteomics, source-protein abundance, flanking sequence, cell-type composition, and antigen-processing features. Tissue information is represented through the task definition and learned sharing structure rather than through direct biological measurements.

Generalization is restricted in several ways. Under the standard split, peptides and parent proteins can recur across tasks and across fitting and evaluation data. Even the peptide-disjoint protocol permits recurrence of parent proteins, studies, platforms, and similar sequence motifs. It is therefore stricter with respect to peptide identity but is not a protein-disjoint, study-disjoint, temporal, external-cohort, or fully independent biological validation. In addition, task-specific output heads mean that both protocols evaluate unseen peptides for previously represented tissue--MHC tasks; they do not test an unseen tissue, an unseen MHC allele, or an unseen tissue--MHC combination.

The human fixed test was used repeatedly during a long model-development process. Although final comparisons follow a frozen protocol, project-level selection may have adapted modeling decisions to this benchmark and may produce optimism not captured by within-run confidence intervals. The mouse benchmark uses a separately frozen internal fixed test, but 81.47\% of its unique test peptides and 88.88\% of its unique parent UniProt identifiers appear in training. It must not be described as strict entity-disjoint or external validation.

MHCflurry and NetMHCpan are frozen general-signal controls, not guaranteed leakage-free competitors. Their pretraining may include IEDB or related immunopeptidomic observations, and the peptide-disjoint split prevents overlap only with the fitting data used by the present tissue-conditioned models. It cannot guarantee that held-out peptides were absent from external pretraining corpora. In addition, the near-tie between MHCflurry and NetMHCpan EL in human is descriptive; without a prespecified direct paired model-comparison test, it does not establish that either external predictor is superior.

Finally, the selected main architectures differ: the TissuePMHC human model is used for human and Factorized MMoE for mouse. Because benchmark scale, MHC system, and data composition also differ, the cross-species results cannot identify species effects. The matched within-species strict controls further show that architecture superiority is not universal: the strongest human auxiliary control is comparable to the full human model, and all three mouse ensembles are comparable.

No shuffled-tissue-label or shuffled-branch control is reported. Such a control requires refitting models after breaking the tissue association and cannot be reconstructed from archived predictions alone. The MHC-only controls answer the narrower completed question of how well the same model families perform when tissue information is absent.

\subsection{Future Work}

Several extensions follow directly from these limitations. First, parent-protein-disjoint, study-disjoint, temporal, and external-cohort splits should progressively test whether performance survives the removal of additional overlap and provenance shortcuts. Shuffled-tissue and shuffled-branch refits should test whether the observed tissue-associated gain exceeds the gain available from task structure alone. Auditable checks against external predictor training corpora would also strengthen the interpretation of frozen general-pMHC controls.

Second, the input representation can be expanded to variable-length peptides and MHC-II ligands. Incorporating tissue transcriptomics, proteomics, source-protein abundance, flanking sequences, proteasomal cleavage features, and cell-type composition could help separate peptide-intrinsic patterns from tissue context. Such additions should be evaluated incrementally under leakage-aware splits so that apparent biological gains are not driven by shared studies or entities.

Third, positive--unlabeled learning and explicit observation models may better reflect the fact that database absence is not biological absence. Calibrated uncertainty estimates could identify tasks and peptides for which the evidence is insufficient. Conditional task representations based on tissue and MHC features, rather than task-specific heads alone, could enable principled evaluation on unseen tissues, alleles, or tissue--MHC combinations.

Finally, controlled cross-species pretraining and adaptation could test whether learned features transfer beyond the shared task construction. Such experiments require harmonized data and matched baselines.
